\documentclass{beamer}

\usepackage{multicol}

\newcommand{\Poincare}{Poincar\'{e}}

\usetheme[progressbar=title]{metropolis}           % Use metropolis theme
\title{The Millennium Prize Problems}
\subtitle{A PGR Seminar in Two Parts}
\date{}
\author{Luke Smallman}
\institute{Cardiff University} \titlegraphic{\hfill\includegraphics[height=1.5cm]{logo}}


\begin{document}
  \maketitle
  \begin{frame}{Structure}
      \begin{multicols}{2}
          \tableofcontents
      \end{multicols}
  \end{frame}

  \section{Introduction}
  \begin{frame}{Introduction}
      The Millennium Prize Problems are problems...
  \end{frame}

  \section{P versus NP}
  \begin{frame}{P vs NP}
      problem statement
  \end{frame}

  \section{Hodge Conjecture}
  \begin{frame}{Hodge Conjecture}
      problem statement
  \end{frame}

  \section{Riemann Hypothesis}
  \begin{frame}{Riemann Hypothesis}
      problem statement
  \end{frame}

  \section{Yang-Mills Existence and Mass Gap}
  \begin{frame}{Yang-Mills Existence and Mass Gap}
      problem statement
  \end{frame}

  \section{Navier-Stokes Existence and Smoothness}
  \begin{frame}{Navier-Stokes Existence and Smoothness}
      problem statement
  \end{frame}

  \section{Birch and Swinnerton-Dyer Conjecture}
  \begin{frame}{Birch and Swinnerton-Dyer Conjecture}
      problem statement
  \end{frame}

  \section{\Poincare{} Conjecture}
  \begin{frame}{\Poincare{} Conjecture}
      problem statement
  \end{frame}


\end{document}
